\documentclass[12pt]{article}
\usepackage[utf8]{inputenc}
\usepackage[brazil]{babel}
\usepackage{amsmath}
\usepackage{amssymb}
\usepackage{graphicx}
\usepackage{float}
\usepackage{geometry}
\geometry{a4paper, margin=25mm}

\begin{document}

% =====================================================================
% CABEÇALHO INSTITUCIONAL
% =====================================================================
\begin{center}
    \begin{minipage}{0.2\textwidth}
        % \includegraphics[width=\textwidth]{logo_ifpb.png}
    \end{minipage}
    \begin{minipage}{0.75\textwidth}
        \begin{center}
            \small\textbf{INSTITUTO FEDERAL DE EDUCAÇÃO, CIÊNCIA E TECNOLOGIA DA PARAÍBA} \\
            \small\textbf{CAMPUS CAMPINA GRANDE} \\
            \small\textbf{CURSO DE ENGENHARIA DE COMPUTAÇÃO} \\
            \small\textbf{DISCIPLINA: CONTROLE E AUTOMAÇÃO I} \\
            \small\textbf{PROFESSOR: PROF. DR. MOACY PEREIRA DA SILVA} \\
        \end{center}
    \end{minipage}
\end{center}

\vspace{0.2cm}
\hrule
\vspace{0.5cm}

% =====================================================================
% TÍTULO DO TRABALHO
% =====================================================================
\begin{center}
    \Large{\textbf{Estudo Dirigido -- Etapa 4: Automação Industrial e Tecnologias}} \\
    \large{\textit{Sistemas Dinâmicos, Modelagem e Aplicações}} \\
    \vspace{0.6cm}

    \large{\textbf{Alunos:} Kevin Ryan e Demétrio de Luna Gurgel} \\
    \large{\textbf{Data:} 11 de Junho de 2026} \\
\end{center}

\vspace{0.8cm}

% =====================================================================
% CONTEÚDO DA ETAPA 4
% =====================================================================

\section{Introdução à Automação Industrial e Elementos Fundamentais}

\subsection{Resumo Teórico sobre Automação Industrial}
A automação industrial consiste na aplicação estratégica de tecnologias de hardware e software para gerenciar, monitorar e controlar processos produtivos de forma autônoma, minimizando a necessidade de intervenção humana direta em tarefas repetitivas ou operacionais. Seus objetivos centrais vão muito além da simples substituição da mão de obra: concentram-se na maximização da produtividade, na garantia de rigorosos padrões de repetibilidade e qualidade, na otimização do uso de insumos energéticos e na mitigação de riscos de acidentes, preservando a integridade física dos operadores em ambientes industriais insalubres ou perigosos.

No paradigma moderno da Indústria 4.0, a automação baseia-se na integração cibernética e em redes de comunicação industriais estruturadas de maneira hierárquica. Isso permite o fluxo contínuo de dados desde o nível de campo até as camadas de tomada de decisão corporativa e gerencial.

\subsection{Análise dos Elementos Fundamentais da Indústria}
A operação coordenada e estável de uma planta automatizada é fortemente dependente da sinergia entre cinco elementos tecnológicos fundamentais, cujas definições e funções são descritas a seguir:

\begin{itemize}
    \item \textbf{Sensores Industriais:} Dispositivos que compõem o nível de campo e atuam como os ``olhos'' e ``ouvidos'' do sistema de automação. Sua função primordial é mensurar grandezas físicas ou químicas do processo (tais como temperatura, pressão, vazão, nível, deslocamento mecânico e presença) e convertê-las em sinais elétricos padronizados (correntes de 4 a 20 mA ou tensões de 0 a 10 V) legíveis pelas unidades de processamento.
    \item \textbf{Controladores Lógicos Programáveis (CLPs):} Dispositivos de computação industrial dedicados, projetados com alta robustez para operar sob condições ambientais severas, resistindo a transientes elétricos, vibrações mecânicas, poeira e elevadas variações térmicas. O CLP executa uma varredura cíclica (\textit{scan}) lendo o estado das variáveis de entrada (sensores), processando interrupções e aplicando algoritmos lógicos ou de controle (como malhas PID discretas) para atualizar os estados das saídas (acionando atuadores, motores e contatores).
    \item \textbf{Sistemas Supervisórios (SCADA):} Softwares de arquitetura centralizada dedicados à Aquisição de Dados e Controle Supervisório. Instalados em servidores ou estações de operação na sala de controle, os sistemas SCADA coletam e tratam os dados compartilhados pelos CLPs em tempo real, fornecendo telas sinóticas para a visualização macroscópica do processo, gerenciamento avançado de alarmes críticos, armazenamento de dados históricos em bancos de dados relacionais e modificações de parâmetros de produção em larga escala.
    \item \textbf{Interfaces Homem-Máquina (IHMs):} Dispositivos gráficos compactos instalados diretamente no chão de fábrica, integrados fisicamente ou próximos aos painéis elétricos das máquinas. Diferente da abrangência global do SCADA, a IHM possui atuação local, servindo como meio de comunicação direta para que os operadores realizem intervenções pontuais, comandos de partida/parada e ajustes manuais específicos daquela célula de trabalho.
    \item \textbf{Protocolos de Comunicação:} Padrões de rede e regras de tráfego de dados indispensáveis para garantir a interoperabilidade entre equipamentos de diferentes fabricantes. No ecossistema industrial de manufatura e processos, destacam-se protocolos determinísticos de campo e de rede como o Modbus (RTU/TCP), Profibus, Profinet, EtherNet/IP e o protocolo MQTT, este último muito empregado no envio direto de métricas de sensores industriais para corretoras em nuvem (IoT Industrial).
\end{itemize}

\section{Linguagens de Programação e Ferramentas para CLPs}

\subsection{Investigação das Linguagens de Programação}
A norma internacional IEC 61131-3 padroniza as linguagens de programação para Controladores Lógicos Programáveis (CLPs), garantindo que engenheiros e técnicos possuam uma base comum de desenvolvimento, independentemente do fabricante do hardware. Entre as linguagens definidas pela norma, destacam-se o \textit{Ladder} e o \textit{Structured Text}:

\begin{itemize}
    \item \textbf{Ladder Diagram (LD):} É uma linguagem gráfica baseada na tradicional lógica de relés eletromecânicos e painéis de contatores. Sua estrutura visual assemelha-se a uma escada (daí o nome \textit{Ladder}), onde as linhas verticais representam a alimentação elétrica e as linhas horizontais (degraus) contêm as condições lógicas (contatos normalmente abertos ou fechados) que acionam as bobinas de saída. Devido à sua natureza visual e intuitiva, é a linguagem mais utilizada na indústria para lógicas de intertravamento, segurança e sequenciamento de máquinas.
    \item \textbf{Structured Text (ST):} É uma linguagem textual de alto nível, com sintaxe estruturada que lembra linguagens de programação tradicionais de engenharia, como C ou Pascal. O ST é extremamente poderoso para aplicações que exigem processamento matemático complexo, manipulação de matrizes, laços de repetição (\textit{FOR, WHILE}) e implementação de malhas de controle avançadas (como algoritmos PID complexos). Para sistemas de controle dinâmico que transcendem a simples lógica de estados booleanos, o ST oferece maior flexibilidade, legibilidade e facilidade de manutenção.
\end{itemize}

\subsection{Ferramentas de Mercado: CodeSys e OpenPLC}
Para o desenvolvimento, simulação e compilação do código antes de enviá-lo ao hardware físico do processo, são utilizados Ambientes de Desenvolvimento Integrado (IDEs) especializados. Analisando o mercado atual e as tendências de pesquisa, destacam-se duas ferramentas:

\begin{itemize}
    \item \textbf{CodeSys (\textit{Controller Development System}):} É uma das ferramentas proprietárias mais consolidadas e robustas do mercado de automação industrial. Atua como um software agnóstico em relação ao hardware, permitindo que o mesmo código desenvolvido nas linguagens da norma IEC 61131-3 seja compilado e executado em CLPs de centenas de fabricantes parceiros. Além da programação lógica, o CodeSys oferece recursos nativos poderosos para simulação virtual da planta, depuração de variáveis em tempo real e criação de telas de IHM em um único ambiente, sendo o padrão em inúmeras aplicações industriais de larga escala e alta confiabilidade.
    \item \textbf{OpenPLC:} Representa um projeto inovador de código aberto (\textit{open-source}) dentro do ecossistema de automação. Criado com o intuito de democratizar o acesso às tecnologias de controle industrial, o OpenPLC permite transformar computadores de propósito geral, placas embarcadas e até microcontroladores em CLPs totalmente funcionais e aderentes às normas da IEC 61131-3. É uma ferramenta de imenso valor para pesquisa acadêmica, desenvolvimento de protótipos em laboratório e integração direta com arquiteturas de \textit{Internet das Coisas Industrial} (IIoT), criando uma ponte direta entre a automação clássica e os sistemas modernos de engenharia de computação.
\end{itemize}

\subsection{Exploração de Aplicação Simulada: Controle de Nível Industrial}

Para relacionar os conceitos teóricos de controle dinâmico com a arquitetura de sistemas industriais reais, propõe-se a análise de uma aplicação simulada: o \textbf{Controle de Nível de um Reservatório Misto}. Este processo é amplamente encontrado em estações de tratamento de água, indústrias químicas e de alimentos.

A dinâmica do sistema envolve o preenchimento de um tanque através de uma bomba centrífuga de rotação variável (variável manipulada) e o escoamento do fluido através de uma válvula de consumo manual (perturbação). O nível atual do líquido é medido por um sensor ultrassônico (variável de processo).

Utilizando um ambiente como o \textit{CodeSys} ou o \textit{OpenPLC}, o desenvolvimento do firmware do controlador é arquitetado explorando o melhor das linguagens da IEC 61131-3:

\begin{itemize}
    \item \textbf{Lógica de Segurança e Intertravamento (Ladder):} A linguagem gráfica \textit{Ladder} é empregada para gerenciar os estados de operação da máquina e proteger o hardware. Através de contatos e bobinas, programa-se o acionamento do modo Automático/Manual, a parada de emergência (que deve cortar imediatamente o sinal da bomba independentemente do controle) e a proteção contra \textit{cavitação} (desligando a bomba caso um sensor de fluxo na tubulação de entrada detecte que ela está operando a seco).
    \item \textbf{Malha de Controle (Structured Text):} A linguagem ST é utilizada para instanciar e processar o bloco de função do controlador PID. O código lê o valor analógico do sensor ultrassônico (escalonado para 0-100\%), calcula o erro em relação ao \textit{setpoint} definido pelo operador e executa as equações de diferença finita (Proporcional, Integral e Derivativa) a cada ciclo de varredura (\textit{scan cycle}) do CLP. O resultado numérico é então repassado para a saída analógica conectada ao inversor de frequência da bomba.
    \item \textbf{Supervisório (IHM Simulada):} Através da ferramenta de visualização integrada do \textit{CodeSys}, projeta-se uma IHM demonstrativa contendo uma representação gráfica do tanque, gráficos de tendência (\textit{trend}) que mostram as curvas da resposta ao degrau da malha fechada, e campos de entrada numéricos para a sintonia fina dos ganhos $K_p$, $K_i$ e $K_d$ em tempo real.
\end{itemize}

Essa abordagem híbrida demonstra como os algoritmos matemáticos de controle estudados previamente são encapsulados e operados dentro de um ecossistema industrial robusto, garantindo não apenas a estabilidade matemática do processo, mas também a integridade operacional do equipamento e a segurança dos operadores.

\section{Tópicos Avançados e Integração de Sistemas Industriais}

\subsection{Temas Avançados em Controle e Automação}
À medida que os processos industriais se tornam mais complexos, dinâmicos e integrados, as estratégias de controle clássicas (como o PID padrão) passam a ser complementadas por abordagens avançadas para lidar com multivariáveis, não-linearidades e grandes volumes de dados. Duas das principais vertentes tecnológicas nesse cenário são o Controle Preditivo Baseado em Modelo (MPC) e a Inteligência Artificial (IA):

\begin{itemize}
    \item \textbf{Controle Preditivo Baseado em Modelo (MPC):} Diferente do PID, que atua de forma reativa com base no erro presente e passado, o MPC opera de forma proativa. Ele utiliza um modelo matemático explícito do processo para prever o comportamento futuro das variáveis de saída ao longo de um horizonte de tempo determinado. A partir dessa previsão, o algoritmo resolve um problema de otimização numérica em tempo real, calculando a sequência ideal de ações de controle que minimizam o desvio em relação às referências, respeitando rigorosamente restrições físicas dos atuadores (como limites de pressão ou válvulas). É amplamente utilizado em indústrias químicas, petroquímicas e de processos multivariáveis complexos.
    \item \textbf{Inteligência Artificial (IA) na Automação:} A IA e o Aprendizado de Máquina (\textit{Machine Learning}) introduzem flexibilidade e capacidade de adaptação aos sistemas de automação. Suas aplicações práticas ramificam-se em frentes fundamentais:
    \begin{itemize}
        \item \textit{Manutenção Preditiva:} Algoritmos de redes neurais e árvores de decisão analisam dados históricos de sensores (vibração, temperatura, corrente) coletados via sistemas SCADA para prever falhas em componentes mecânicos antes que elas ocorram, reduzindo paradas não programadas.
        \item \textit{Visão Computacional:} Redes neurais convolucionais (CNNs) integradas a câmeras no chão de fábrica realizam o controle de qualidade automatizado em alta velocidade, detectando imperfeições estéticas ou dimensionais em produtos diretamente na esteira de produção.
        \item \textit{Controle Inteligente (Edge AI):} Redes neurais embarcadas em hardware local (como microcontroladores robustos) podem atuar na otimização de parâmetros de sintonia de controladores industriais em tempo real, adaptando os ganhos do sistema à medida que os equipamentos sofrem desgaste físico.
    \end{itemize}
\end{itemize}

\subsection{Análise Final: A Integração no Ecossistema Industrial Moderno}
A evolução do aprendizado ao longo deste estudo dirigido permite consolidar uma visão sistêmica sobre a engenharia de controle e automação. O desenvolvimento de uma solução de engenharia de excelência não ocorre de forma isolada, mas sim através de uma cadeia integrada de conceitos:

A jornada inicia-se com a \textbf{Modelagem Matemática}, onde as leis da física e equações diferenciais capturam a essência dinâmica e as limitações do sistema real. Sem um modelo preciso, o projeto de qualquer estratégia de controle torna-se puramente empírico e arriscado.

Em seguida, o \textbf{Controle (PID/MPC)} utiliza essa base matemática para projetar a inteligência do sistema, calculando como mitigar erros, estabilizar dinâmicas oscilatórias e garantir que o processo atinja os objetivos de operação de forma rápida e segura.

Por fim, a \textbf{Automação Industrial} fornece a infraestrutura tecnológica necessária para traduzir a matemática e os algoritmos em realidade física. É através de redes de comunicação determinísticas, sensores confiáveis, IHMs locais e o processamento em tempo real de CLPs robustos que as leis de controle ganham vida no chão de fábrica de forma segura, escalável e alinhada com as demandas de eficiência da indústria moderna.

\end{document}